% Section IV: type rule, learned per-sample routing, and conditional analysis.
\section{Question-Conditioned Evidence Routing}\label{sec:routing}

The two branches in Section~\ref{sec:system} differ in both the
information they transmit and the computation they require. Evidence
routing selects a branch before transmission. When detection evidence
is sufficient to answer the question, the receiver can avoid VLM
inference. We first introduce a question-type rule that requires no
router training. We then develop a per-sample router that uses the
current detector output to refine branch selection. The final subsection
analyzes the conditions under which evidence selection can reduce
energy without reducing accuracy.

\subsection{Question-Type-Based Evidence Selection}\label{subsec:rule}

Counting, comparison, co-presence, and threshold questions can be
expressed as functions of object counts. The detection-evidence branch
provides these counts to a symbolic decoder and avoids image-based
inference at the receiver. Presence questions can also be answered
from detections. However, a failure to detect an object does not mean
that the object is absent from the scene. The image branch retains
visual information that the detector may have missed. This difference
motivates branch selection based on question type.

Let $t(q)$ denote the type of question $q$. Let
$\mathcal T_{\mathrm{sym}}$ contain counting, comparison, co-presence,
and threshold, and let $\mathcal T_{\mathrm{pres}}$ contain presence.
The type rule is
\begin{equation}
r(t(q))=
\begin{cases}
\mathrm{det}, & t(q)\in\mathcal T_{\mathrm{sym}},\\
\mathrm{img}, & t(q)\in\mathcal T_{\mathrm{pres}}.
\end{cases}
\label{eq:rule}
\end{equation}
The rule uses neither channel-state information (CSI) nor receiver
feedback. If it selects the detection-evidence branch, the UAV runs
the detector and transmits the resulting packet. Otherwise, it
transmits the image without running the detector for routing. Thus,
the decision also determines whether onboard detection is required.
Equation~\eqref{eq:energy} includes this detection cost.

Equation~\eqref{eq:rule} defines a fixed rule. It does not guarantee
that the selected branch is more accurate for every question of a
given type. Missed objects, false detections, and image distortion can
change which evidence is more useful for an individual query.
For threshold questions, detection errors can move the estimated count
across the threshold specified in the question and change the answer. A type-only
rule cannot account for these sample-specific differences. This
motivates the learned router below.

\subsection{Per-Sample Evidence Routing}\label{subsec:persample}

The per-sample router predicts whether each branch will answer correctly
using information available before transmission. For each query, the
UAV first computes $Z=g(I)$ and constructs a feature vector $x$.
The feature vector has $18$ entries: five question-type indicators,
ten target-class indicators, the assumed SNR in dB divided by $20$,
the target-class detection count clipped at $60$ and divided by $60$,
and an indicator of whether that count is nonzero. The count is
computed from the original detector records $Z$, before channel
impairment. It is not the count recovered at the receiver.
Ground-truth answers, annotation-derived scene attributes, and
receiver outputs are excluded. The assumed SNR is a transmitter-side
input; no receiver VLM pass is required before selection.

For each branch $s\in\mathcal S=\{\mathrm{det},\mathrm{img}\}$, let
$\widehat p_s(x)$ estimate the probability that its answer is correct.
Training uses paired binary labels $y_{i,s}\in\{0,1\}$ from the
recorded answers of both branches. Each pair corresponds to the same
image, question, and channel condition. Each label indicates whether
the answer meets the scoring rule in Section~\ref{subsec:setup}.
The labels describe the correctness of each branch, rather than
identify one preferred branch. Both answers can therefore be correct
or incorrect.

We train the correctness predictors using binary cross-entropy.
For $N_{\mathrm{tr}}$ training decisions, the data-fitting term is
\begin{equation}
\mathcal L_{\mathrm{BCE}}
=-\frac{1}{2N_{\mathrm{tr}}}
\sum_{i=1}^{N_{\mathrm{tr}}}\sum_{s\in\mathcal S}
\left[y_{i,s}\log\widehat p_s(x_i)
+(1-y_{i,s})\log(1-\widehat p_s(x_i))\right].
\label{eq:router-loss}
\end{equation}
We use compact multilayer perceptron (MLP) predictors.
Each predictor has two hidden layers with $32$ and $16$ units and
uses ReLU activations. We train the predictors with Adam.
Predictor fitting uses training images, and early stopping uses a
separate set of validation images. All questions and SNR replays of
an image stay in the same partition. Section~\ref{subsec:setup}
specifies the split, count calibration, and repeated runs.
Predicting branch performance before transmission is related to
predictive adaptation in image semantic communication~\cite{zhang2023padc}.
Here, the prediction target is answer correctness rather than image
reconstruction quality.

The default decision selects the branch with the higher predicted
probability of a correct answer:
\begin{equation}
\widehat s(x)=\arg\max_{s\in\mathcal S}\widehat p_s(x).
\label{eq:argmax}
\end{equation}
An energy-aware extension uses the same predictors and includes a
weighted energy cost in the decision:
\begin{equation}
\widehat s_\lambda(x)=
\arg\max_{s\in\mathcal S}
\left[\widehat p_s(x)-\lambda E_\rho(s,\gamma)\right],
\qquad \lambda\ge0,
\label{eq:argmax-lambda}
\end{equation}
where $\lambda$ has units J$^{-1}$ and $E_\rho(s,\gamma)$ is the
energy for selector $\rho$ in \eqref{eq:energy}. Setting $\lambda=0$
gives \eqref{eq:argmax}. We report this setting alongside
energy-aware operating points. A positive energy price favors the lower-cost branch when
its predicted accuracy is sufficiently close to that of the other
branch. Section~\ref{subsec:energyctl} evaluates a fixed price grid
and a validation-selected operating point. Costs are evaluated
per query; their image dependence is suppressed in the notation.
The rule does not impose a hard per-query budget or guarantee
a test-set accuracy loss.

Algorithm~\ref{alg:selector} summarizes the two routing options.
Unlike the type rule, per-sample routing runs the detector for every
query, even when it selects the image branch. This detection cost is
the same for both branch choices and does not change their score
difference. It must still be included in the total energy.
The selected branch then follows the transmission and
answer-generation models of Section~\ref{sec:system}. Training uses
paired outcomes, but inference transmits only the selected evidence.
The experiments evaluate this procedure by replaying recorded branch
outcomes rather than running a live control loop. The type-level
lookup table and linear predictor serve as baselines.
Section~\ref{subsec:setup} defines these baselines.

\begin{algorithm}[t]
\caption{Question-conditioned evidence routing for one query}
\label{alg:selector}
\begin{algorithmic}[1]
\REQUIRE image $I$; question $q$ and type $t(q)$; routing mode
$\rho\in\{\mathrm{rule},\mathrm{MLP}\}$;
for MLP, target class, assumed SNR $\gamma$, predictors
$\widehat p_s$, and energy price $\lambda$ (default $0$)
\IF{$\rho=\mathrm{rule}$}
  \STATE $\widehat s\leftarrow r(t(q))$ using \eqref{eq:rule}
\ELSE
  \STATE compute $Z=g(I)$ and construct features $x$
  \STATE predict $\widehat p_{\mathrm{det}}(x)$ and $\widehat p_{\mathrm{img}}(x)$
  \STATE select $\widehat s$ using \eqref{eq:argmax-lambda}
\ENDIF
\IF{$\widehat s=\mathrm{det}$}
  \IF{$\rho=\mathrm{rule}$}
    \STATE compute $Z=g(I)$
  \ENDIF
  \STATE transmit $Z$; answer from received evidence $\widetilde Z$
\ELSE
  \STATE transmit the rate-adaptive coded image; answer from $\widehat I$
\ENDIF
\end{algorithmic}
\end{algorithm}

\subsection{Analysis of Evidence Selection}\label{sec:theory}

\emph{Information required by the question.}
Consider a counting question with ground truth $A_q=N_c(I)$.
If detection recovers the true count and transmission preserves the
relevant records, then
$\widehat N_c(\widetilde Z)=N_c(I)$, and
\begin{equation}
H\!\left(A_q\mid\widehat N_c(\widetilde Z)\right)=0.
\label{eq:count-sufficiency}
\end{equation}
Thus, under these ideal conditions, the detection evidence contains
all information needed for that answer. The image is not required.
The same result applies to comparison, co-presence, and
threshold questions when all relevant class counts are recovered.
This result states when detection evidence is sufficient. It does
not guarantee that a practical detector or a noisy link meets these
conditions.

In practice, the detection-evidence branch is affected by detection
errors and corruption or loss of transmitted records. The image
branch is affected by compression, channel degradation, and VLM
answering errors. For presence questions, an object missed by the
detector may be absent from the packet but still visible in the
image. This explains a potential advantage of
image transmission. It does not prove that the image branch is
always more accurate. The comparison depends on false positives, missed
objects, the channel, and the actual receiver. Section~\ref{subsec:comp}
therefore measures branch accuracy by question type instead of
inferring it from detector recall alone.

\emph{Conditions for energy reduction.}
\label{subsec:pareto-theory}
Fix a channel condition and the computation model. Let question type
$t\in\mathcal T$ have probability $w_t>0$, with
$\sum_{t\in\mathcal T}w_t=1$. Let $a_s(t)$ and $e_s(t)$ denote the
expected accuracy and per-query energy of branch $s$, respectively.
Here the costs correspond to type-only selection: detection is
performed only when its branch is selected. A deterministic
type-based policy $\pi:\mathcal T\to\mathcal S$ has
\begin{equation}
A(\pi)=\sum_{t\in\mathcal T}w_t a_{\pi(t)}(t),
\qquad
E(\pi)=\sum_{t\in\mathcal T}w_t e_{\pi(t)}(t).
\label{eq:policy}
\end{equation}

\begin{proposition}[Conditional energy--accuracy dominance]\label{prop:pareto}
Let $\mathcal T=\mathcal T_{\mathrm{sym}}\cup\mathcal T_{\mathrm{pres}}$,
where the two sets in \eqref{eq:rule} are disjoint and nonempty. Assume that
$e_{\mathrm{det}}(t)<e_{\mathrm{img}}(t)$ for every $t$, and that
\begin{equation}
\begin{aligned}
a_{\mathrm{det}}(t)&\ge a_{\mathrm{img}}(t),
&&t\in\mathcal T_{\mathrm{sym}},\\
a_{\mathrm{img}}(t)&>a_{\mathrm{det}}(t),
&&t\in\mathcal T_{\mathrm{pres}}.
\end{aligned}
\label{eq:ordering}
\end{equation}
Then the type rule $r$ maximizes $A(\pi)$ over deterministic type-based
policies. Among the policies with maximum accuracy, $r$ has the lowest
energy. Relative to the fixed-image policy $\pi_{\mathrm{img}}$,
it satisfies $A(r)\ge A(\pi_{\mathrm{img}})$ and
$E(r)<E(\pi_{\mathrm{img}})$.
\end{proposition}

\emph{Proof.}
Each term in the accuracy sum in \eqref{eq:policy} can be maximized
separately. Under \eqref{eq:ordering}, $r$ selects a branch with maximum
accuracy for each type. When the two branches have equal accuracy,
the detection-evidence branch has lower energy. On the symbolic
types, $r$ replaces image transmission with detection evidence without
reducing accuracy. These types have positive total probability, so
this replacement strictly reduces the expected energy.
These observations prove the stated results. \hfill$\square$

The proposition states sufficient conditions, not a general
optimality guarantee for \eqref{eq:rule}. In particular, the accuracy
condition must also hold for threshold questions, which are part of
$\mathcal T_{\mathrm{sym}}$. A small observed accuracy difference does
not establish equality or confirm this condition. The
experiments assess the rule and the learned router directly under
the stated energy model. The MLP uses information beyond question
type and has different computation costs. This proposition therefore
does not prove the gains measured for the MLP.
